\section{广义不等式}

\begin{BoxDefinition}[正常锥]
    若锥$K\subseteq\vb{R}^n$满足以下条件，则称为正常锥（Proper Cone）
    \begin{itemize}
        \item $K$是凸的（Convex）。
        \item $K$是闭的（Closed），即集合是包含自身边界的。
        \item $K$是实心的（Solid），即集合内没有空洞。
        \item $K$是尖的（Pointed），即集合内不包含贯通的直线。
    \end{itemize}
\end{BoxDefinition}

通常来说，我们能想象的凸锥的性质都已经很好了，都是正常锥！需要解释的是为什么尖和集合内不存在贯通直线有关：当一个锥的角度越来越大，达到$180^\circ$时，锥角就不再是尖锐的了。

集合上如何比较大小是一个麻烦的问题：我们可以说$5>3$，但很难定义$(5,1)>(3,1)$。正常锥的意义是，我们可以定义偏序（Partial Orderning）的广义不等式（Generalized Inequality）！

\begin{BoxDefinition}[广义不等式]
    对于正常锥$K\subseteq\R$，定义广义不等式$\succeq_K$
    \begin{Equation}[广义不等式]
        \vb{y}\succeq_K\vb{x}\quad\Longleftrightarrow\quad\vb{y}-\vb{x}\in K
    \end{Equation}
\end{BoxDefinition}

比较常用的定义广义不等式的正常锥有

\begin{itemize}
    \item 适用于向量的第一象限锥$\R^n_{+}$上定义的$\succeq_{\R^n_{+}}$不等式（等价于逐元素）。
    \item 适用于矩阵的半正定锥$\Symm^n_{+}$上定义的$\succeq_{\Symm^n_{+}}$不等式（等价于二次型）。
\end{itemize}

偏序的含义是，我们可以比较两个元素，但我们会发现某些元素无论是$a\succeq b$还是$b\succeq a$都不成立，因此无法线性排序所有元素！例如$(2,1)\succeq_{\R^n_{+}}(1,2)$和$(1,2)\succeq_{\R^{n}_{+}}(2,1)$就都不成立。

\begin{BoxProperty}[广义不等式的可加性]
    $\succeq_{K}$满足可加性：若$\vb{x}\succeq_{K}\vb{y}$且$\vb{u}\succeq_{K}\vb{v}$，那么有$\vb{x}+\vb{u}\succeq \vb{y}+\vb{v}$成立。
\end{BoxProperty}

\begin{BoxProperty}[广义不等式的传递性]
    $\succeq_{K}$满足传递性：若$\vb{x}\succeq_{K}\vb{y}$且$\vb{y}\succeq_{K}\vb{z}$，那么有$\vb{x}\succeq \vb{z}$成立。
\end{BoxProperty}

\begin{BoxProperty}[广义不等式的放缩性]
    $\succeq_{K}$满足放缩性：若$\vb{x}\succeq_{K}\vb{y}$且$\alpha\geq 0$，那么有$\alpha \vb{x}\succeq \alpha \vb{z}$成立。
\end{BoxProperty}

\begin{BoxProperty}[广义不等式的反称性]
    $\succeq_{K}$满足反称性：若$\vb{x}\succeq_{K}\vb{y}$且$\vb{y}\succeq_{K}\vb{x}$，那么有$\vb{x}=\vb{y}$成立。
\end{BoxProperty}

\begin{BoxProperty}[广义不等式的自反性]
    $\succeq_{K}$满足自反性：始终有$\vb{x}\succeq_{K}\vb{x}$成立。
\end{BoxProperty}
